
\documentclass{article}
\usepackage{colortbl}
\usepackage{makecell}
\usepackage{multirow}
\usepackage{supertabular}

\begin{document}

\newcounter{utterance}

\centering \large Interaction Transcript for game `dond', experiment `coop\_en', episode 10 with Qwen3{-}32B{-}t0.0.
\vspace{24pt}

{ \footnotesize  \setcounter{utterance}{1}
\setlength{\tabcolsep}{0pt}
\begin{supertabular}{c@{$\;$}|p{.15\linewidth}@{}p{.15\linewidth}p{.15\linewidth}p{.15\linewidth}p{.15\linewidth}p{.15\linewidth}}
    \# & $\;$A & \multicolumn{4}{c}{Game Master} & $\;\:$B\\
    \hline

    \theutterance \stepcounter{utterance}  
    & & \multicolumn{4}{p{0.6\linewidth}}{
        \cellcolor[rgb]{0.9,0.9,0.9}{
            \makecell[{{p{\linewidth}}}]{
                \texttt{\tiny{[P1$\langle$GM]}}
                \texttt{You are playing a cooperative negotiation game in which you have to agree on how to divide a set of items among you and another player.} \\
\\ 
\texttt{Rules:} \\
\texttt{(a) You and the other player are given a set of items. Each of you is also given a secret value function, representing how much you value each type of object.} \\
\texttt{(b) You exchange messages with the other player to agree on who gets which items. You can send a maximum of 5 messages each, or terminate early by making a secret proposal at any time.} \\
\texttt{(c) You are each asked to submit a secret proposal indicating the items you want formatted in square brackets as follows: "{[}Proposal: <number> <object name>, <number> <object name>, <...>{]}"} \\
\texttt{(d) If your proposals are complementary, i.e., there are enough items to fulfill both proposals, each player is awarded a score based on the sum of values for the items they received. Otherwise, both of you get zero points.} \\
\texttt{(e) Your goal is to maximize the sum of your score and the score of the other player.} \\
\\ 
\texttt{Let us start.} \\
\\ 
\texttt{The set of available items is:} \\
\\ 
\texttt{2 pliers, 1 chalk, 1 nail, 1 bucket} \\
\\ 
\texttt{Your secret value function is:} \\
\\ 
\texttt{plier: 2, chalk: 3, nail: 2, bucket: 1} \\
\\ 
\texttt{IMPORTANT: Your messages, unless it is the secret proposal, are directly transmitted to the other player, so do not include any response to the rules or text announcing your message. To make a secret proposal, use the indicated format. Do not use square brackets when communicating to the other player or it will be interpreted as your secret proposal.} \\
            }
        }
    }
    & \\ \\

    \theutterance \stepcounter{utterance}  
    & \multicolumn{4}{p{0.6\linewidth}}{
        \cellcolor[rgb]{0.9,0.9,0.9}{
            \makecell[{{p{\linewidth}}}]{
                \texttt{\tiny{[P1$\rangle$GM]}}
                \texttt{<think>} \\
\texttt{Okay, let's see. I need to divide these items between me and another player. The items are 2 pliers, 1 chalk, 1 nail, and 1 bucket. My value function is plier: 2, chalk: 3, nail: 2, bucket: 1. The other player's values are secret, so I don't know what they value. The goal is to maximize the total score for both of us, which means we need to find a division where our proposals are complementary.} \\
\\ 
\texttt{First, I should figure out what the other player might value. Since I don't have their value function, I need to make some assumptions or ask for information. But since I can only send 5 messages, I need to be efficient. Maybe I can ask them what they value most, but they might not tell me the truth. Alternatively, I can propose a division and see if they accept.} \\
\\ 
\texttt{Looking at my own values, the highest value is chalk at 3, then pliers and nail at 2 each, and bucket at 1. So I might prefer getting the chalk and maybe one plier. But if the other player also values chalk, we might have a conflict. Alternatively, if they value the bucket or nail more, we can split accordingly.} \\
\\ 
\texttt{Since there are two pliers, maybe we can each take one. But if the other player doesn't value pliers, they might prefer other items. Let me think about possible splits. For example, if I take the chalk and one plier, that's 3 + 2 = 5. Then the other player could take the other plier, nail, and bucket. Their total would depend on their values. But if their values are such that they value the nail and bucket more, that could work. Alternatively, if they value the pliers more, maybe they want both, but there are two, so we can split them.} \\
\\ 
\texttt{Wait, the key is that our proposals need to be complementary. So if I propose taking certain items, the other player's proposal must not overlap. So I need to make sure that the items I want don't conflict with theirs. Since I don't know their values, maybe I should propose a division that allows flexibility. For example, if I take the chalk (which is unique), and one plier, then the other player can take the remaining plier, nail, and bucket. But if they also} \\
            }
        }
    }
    & & \\ \\

    \theutterance \stepcounter{utterance}  
    & & & \multicolumn{2}{p{0.3\linewidth}}{
        \cellcolor[rgb]{0.9,0.9,0.9}{
            \makecell[{{p{\linewidth}}}]{
                \texttt{\tiny{[GM$|$GM]}}
                \texttt{continue} \\
            }
        }
    }
    & & \\ \\

    \theutterance \stepcounter{utterance}  
    & & & \multicolumn{2}{p{0.3\linewidth}}{
        \cellcolor[rgb]{0.9,0.9,0.9}{
            \makecell[{{p{\linewidth}}}]{
                \texttt{\tiny{[GM$|$GM]}}
                \texttt{<think>} \\
\texttt{Okay, let's see. I need to divide these items between me and another player. The items are 2 pliers, 1 chalk, 1 nail, and 1 bucket. My value function is plier: 2, chalk: 3, nail: 2, bucket: 1. The other player's values are secret, so I don't know what they value. The goal is to maximize the total score for both of us, which means we need to find a division where our proposals are complementary.} \\
\\ 
\texttt{First, I should figure out what the other player might value. Since I don't have their value function, I need to make some assumptions or ask for information. But since I can only send 5 messages, I need to be efficient. Maybe I can ask them what they value most, but they might not tell me the truth. Alternatively, I can propose a division and see if they accept.} \\
\\ 
\texttt{Looking at my own values, the highest value is chalk at 3, then pliers and nail at 2 each, and bucket at 1. So I might prefer getting the chalk and maybe one plier. But if the other player also values chalk, we might have a conflict. Alternatively, if they value the bucket or nail more, we can split accordingly.} \\
\\ 
\texttt{Since there are two pliers, maybe we can each take one. But if the other player doesn't value pliers, they might prefer other items. Let me think about possible splits. For example, if I take the chalk and one plier, that's 3 + 2 = 5. Then the other player could take the other plier, nail, and bucket. Their total would depend on their values. But if their values are such that they value the nail and bucket more, that could work. Alternatively, if they value the pliers more, maybe they want both, but there are two, so we can split them.} \\
\\ 
\texttt{Wait, the key is that our proposals need to be complementary. So if I propose taking certain items, the other player's proposal must not overlap. So I need to make sure that the items I want don't conflict with theirs. Since I don't know their values, maybe I should propose a division that allows flexibility. For example, if I take the chalk (which is unique), and one plier, then the other player can take the remaining plier, nail, and bucket. But if they also} \\
            }
        }
    }
    & & \\ \\

    \theutterance \stepcounter{utterance}  
    & & \multicolumn{4}{p{0.6\linewidth}}{
        \cellcolor[rgb]{0.9,0.9,0.9}{
            \makecell[{{p{\linewidth}}}]{
                \texttt{\tiny{[GM$\rangle$P2]}}
                \texttt{You are playing a cooperative negotiation game in which you have to agree on how to divide a set of items among you and another player.} \\
\\ 
\texttt{Rules:} \\
\texttt{(a) You and the other player are given a set of items. Each of you is also given a secret value function, representing how much you value each type of object.} \\
\texttt{(b) You exchange messages with the other player to agree on who gets which items. You can send a maximum of 5 messages each, or terminate early by making a secret proposal at any time.} \\
\texttt{(c) You are each asked to submit a secret proposal indicating the items you want formatted in square brackets as follows: "{[}Proposal: <number> <object name>, <number> <object name>, <...>{]}"} \\
\texttt{(d) If your proposals are complementary, i.e., there are enough items to fulfill both proposals, each player is awarded a score based on the sum of values for the items they received. Otherwise, both of you get zero points.} \\
\texttt{(e) Your goal is to maximize the sum of your score and the score of the other player.} \\
\\ 
\texttt{Let us start.} \\
\\ 
\texttt{The set of available items is:} \\
\\ 
\texttt{2 pliers, 1 chalk, 1 nail, 1 bucket} \\
\\ 
\texttt{Your secret value function is:} \\
\\ 
\texttt{plier: 0, chalk: 2, nail: 5, bucket: 3} \\
\\ 
\texttt{IMPORTANT: Your messages, unless it is the secret proposal, are directly transmitted to the other player, so do not include any response to the rules or text announcing your message. To make a secret proposal, use the indicated format. Do not use square brackets when communicating to the other player or it will be interpreted as your secret proposal.} \\
\\ 
\\ 
\texttt{<think>} \\
\texttt{Okay, let's see. I need to divide these items between me and another player. The items are 2 pliers, 1 chalk, 1 nail, and 1 bucket. My value function is plier: 2, chalk: 3, nail: 2, bucket: 1. The other player's values are secret, so I don't know what they value. The goal is to maximize the total score for both of us, which means we need to find a division where our proposals are complementary.} \\
\\ 
\texttt{First, I should figure out what the other player might value. Since I don't have their value function, I need to make some assumptions or ask for information. But since I can only send 5 messages, I need to be efficient. Maybe I can ask them what they value most, but they might not tell me the truth. Alternatively, I can propose a division and see if they accept.} \\
\\ 
\texttt{Looking at my own values, the highest value is chalk at 3, then pliers and nail at 2 each, and bucket at 1. So I might prefer getting the chalk and maybe one plier. But if the other player also values chalk, we might have a conflict. Alternatively, if they value the bucket or nail more, we can split accordingly.} \\
\\ 
\texttt{Since there are two pliers, maybe we can each take one. But if the other player doesn't value pliers, they might prefer other items. Let me think about possible splits. For example, if I take the chalk and one plier, that's 3 + 2 = 5. Then the other player could take the other plier, nail, and bucket. Their total would depend on their values. But if their values are such that they value the nail and bucket more, that could work. Alternatively, if they value the pliers more, maybe they want both, but there are two, so we can split them.} \\
\\ 
\texttt{Wait, the key is that our proposals need to be complementary. So if I propose taking certain items, the other player's proposal must not overlap. So I need to make sure that the items I want don't conflict with theirs. Since I don't know their values, maybe I should propose a division that allows flexibility. For example, if I take the chalk (which is unique), and one plier, then the other player can take the remaining plier, nail, and bucket. But if they also} \\
            }
        }
    }
    & \\ \\

    \theutterance \stepcounter{utterance}  
    & & & \multicolumn{4}{p{0.6\linewidth}}{
        \cellcolor[rgb]{0.9,0.9,0.9}{
            \makecell[{{p{\linewidth}}}]{
                \texttt{\tiny{[GM$\langle$P2]}}
                \texttt{<think>} \\
\texttt{Okay, let's start by understanding the problem. I need to divide 2 pliers, 1 chalk, 1 nail, and 1 bucket between me and another player. My values are plier: 2, chalk: 3, nail: 2, bucket: 1. The other player's values are secret. The goal is to maximize the total score for both of us, which means our proposals must be complementary.} \\
\\ 
\texttt{First, I should consider what items I value the most. Chalk is highest at 3, then pliers and nail at 2 each, and bucket at 1. Since there's only one chalk, if I take it, the other player can't. But maybe the other player values chalk more. However, since I don't know their values, I need to make a proposal that allows for a complementary split.} \\
\\ 
\texttt{The items are limited: two pliers, and one each of the others. If I take one plier, the other can take the other. That way, we don't conflict on pliers. Then, the other items (chalk, nail, bucket) are unique. If I take chalk, the other can take nail and bucket. Let's calculate the total value in that case. My value would be chalk (3) + one plier (2) = 5. The other player's value would depend on their values for nail and bucket. If they value nail and bucket highly, that's good. But since I don't know their values, maybe I should propose taking chalk and one plier, and let them take the rest. Alternatively, maybe I should take the nail and bucket, but that might not be optimal for me.} \\
\\ 
\texttt{Another approach: since the total score is the sum of both players' scores, I should aim for a division where the combined value is maximized. For example, if I take chalk (3) and one plier (2), and the other takes the other plier, nail, and bucket. The total would be 3+2 (me) plus their values for plier + nail + bucket. But without knowing their values, it's hard to calculate. However, since the other player's values are secret, maybe the best I can do is propose a division that allows them to take items they might value more, while I take what I value most.} \\
\\ 
\texttt{Alternatively, maybe I should propose taking the chalk and the nail, which are} \\
            }
        }
    }
     \\ \\

    \theutterance \stepcounter{utterance}  
    & & & \multicolumn{2}{p{0.3\linewidth}}{
        \cellcolor[rgb]{0.9,0.9,0.9}{
            \makecell[{{p{\linewidth}}}]{
                \texttt{\tiny{[GM$|$GM]}}
                \texttt{continue} \\
            }
        }
    }
    & & \\ \\

    \theutterance \stepcounter{utterance}  
    & & & \multicolumn{2}{p{0.3\linewidth}}{
        \cellcolor[rgb]{0.9,0.9,0.9}{
            \makecell[{{p{\linewidth}}}]{
                \texttt{\tiny{[GM$|$GM]}}
                \texttt{<think>} \\
\texttt{Okay, let's start by understanding the problem. I need to divide 2 pliers, 1 chalk, 1 nail, and 1 bucket between me and another player. My values are plier: 2, chalk: 3, nail: 2, bucket: 1. The other player's values are secret. The goal is to maximize the total score for both of us, which means our proposals must be complementary.} \\
\\ 
\texttt{First, I should consider what items I value the most. Chalk is highest at 3, then pliers and nail at 2 each, and bucket at 1. Since there's only one chalk, if I take it, the other player can't. But maybe the other player values chalk more. However, since I don't know their values, I need to make a proposal that allows for a complementary split.} \\
\\ 
\texttt{The items are limited: two pliers, and one each of the others. If I take one plier, the other can take the other. That way, we don't conflict on pliers. Then, the other items (chalk, nail, bucket) are unique. If I take chalk, the other can take nail and bucket. Let's calculate the total value in that case. My value would be chalk (3) + one plier (2) = 5. The other player's value would depend on their values for nail and bucket. If they value nail and bucket highly, that's good. But since I don't know their values, maybe I should propose taking chalk and one plier, and let them take the rest. Alternatively, maybe I should take the nail and bucket, but that might not be optimal for me.} \\
\\ 
\texttt{Another approach: since the total score is the sum of both players' scores, I should aim for a division where the combined value is maximized. For example, if I take chalk (3) and one plier (2), and the other takes the other plier, nail, and bucket. The total would be 3+2 (me) plus their values for plier + nail + bucket. But without knowing their values, it's hard to calculate. However, since the other player's values are secret, maybe the best I can do is propose a division that allows them to take items they might value more, while I take what I value most.} \\
\\ 
\texttt{Alternatively, maybe I should propose taking the chalk and the nail, which are} \\
            }
        }
    }
    & & \\ \\

    \theutterance \stepcounter{utterance}  
    & & \multicolumn{4}{p{0.6\linewidth}}{
        \cellcolor[rgb]{0.9,0.9,0.9}{
            \makecell[{{p{\linewidth}}}]{
                \texttt{\tiny{[P1$\langle$GM]}}
                \texttt{<think>} \\
\texttt{Okay, let's start by understanding the problem. I need to divide 2 pliers, 1 chalk, 1 nail, and 1 bucket between me and another player. My values are plier: 2, chalk: 3, nail: 2, bucket: 1. The other player's values are secret. The goal is to maximize the total score for both of us, which means our proposals must be complementary.} \\
\\ 
\texttt{First, I should consider what items I value the most. Chalk is highest at 3, then pliers and nail at 2 each, and bucket at 1. Since there's only one chalk, if I take it, the other player can't. But maybe the other player values chalk more. However, since I don't know their values, I need to make a proposal that allows for a complementary split.} \\
\\ 
\texttt{The items are limited: two pliers, and one each of the others. If I take one plier, the other can take the other. That way, we don't conflict on pliers. Then, the other items (chalk, nail, bucket) are unique. If I take chalk, the other can take nail and bucket. Let's calculate the total value in that case. My value would be chalk (3) + one plier (2) = 5. The other player's value would depend on their values for nail and bucket. If they value nail and bucket highly, that's good. But since I don't know their values, maybe I should propose taking chalk and one plier, and let them take the rest. Alternatively, maybe I should take the nail and bucket, but that might not be optimal for me.} \\
\\ 
\texttt{Another approach: since the total score is the sum of both players' scores, I should aim for a division where the combined value is maximized. For example, if I take chalk (3) and one plier (2), and the other takes the other plier, nail, and bucket. The total would be 3+2 (me) plus their values for plier + nail + bucket. But without knowing their values, it's hard to calculate. However, since the other player's values are secret, maybe the best I can do is propose a division that allows them to take items they might value more, while I take what I value most.} \\
\\ 
\texttt{Alternatively, maybe I should propose taking the chalk and the nail, which are} \\
            }
        }
    }
    & \\ \\

    \theutterance \stepcounter{utterance}  
    & \multicolumn{4}{p{0.6\linewidth}}{
        \cellcolor[rgb]{0.9,0.9,0.9}{
            \makecell[{{p{\linewidth}}}]{
                \texttt{\tiny{[P1$\rangle$GM]}}
                \texttt{<think>} \\
\texttt{Okay, let's try to figure this out. I need to divide the items between me and the other player. My values are plier:2, chalk:3, nail:2, bucket:1. The other player's values are secret. The key is that our proposals must be complementary, so we can't both want the same items. My goal is to maximize the total score for both of us.} \\
\\ 
\texttt{First, I should consider which items I value the most. Chalk is highest at 3, then pliers and nail at 2 each, and bucket at 1. Since there's only one chalk, if I take it, the other player can't. But maybe the other player values chalk more than I do. However, since I don't know their values, I have to make an educated guess.} \\
\\ 
\texttt{The items are 2 pliers, 1 chalk, 1 nail, 1 bucket. Let's think about possible splits. If I take the chalk (3) and one plier (2), that's 5 points for me. Then the other player could take the remaining plier, nail, and bucket. Their total would depend on their values for those items. If they value the plier, nail, and bucket highly, that could work. But if they don't, maybe they'd prefer something else.} \\
\\ 
\texttt{Alternatively, if I take the nail (2) and bucket (1), that's 3 points, but maybe the other player can take the chalk and pliers. But that might not be optimal for me. Wait, the total combined value would be higher if I take the chalk and plier, and they take the other plier, nail, and bucket. Let's see: my 3+2=5, and their plier (2) + nail (2) + bucket (1) = 5. Total 10. If I took nail and bucket (3) and they took chalk and pliers (3+2=5), total is 8. So the first option is better for total.} \\
\\ 
\texttt{But since I don't know their values, maybe they value the nail or bucket more. But without knowing, I have to assume that the other player's values are such that taking the remaining items after I take chalk and one plier would give them a good score. So maybe I should propose taking chalk and one plier. Then the other player can take the other} \\
            }
        }
    }
    & & \\ \\

    \theutterance \stepcounter{utterance}  
    & & & \multicolumn{2}{p{0.3\linewidth}}{
        \cellcolor[rgb]{0.9,0.9,0.9}{
            \makecell[{{p{\linewidth}}}]{
                \texttt{\tiny{[GM$|$GM]}}
                \texttt{continue} \\
            }
        }
    }
    & & \\ \\

    \theutterance \stepcounter{utterance}  
    & & & \multicolumn{2}{p{0.3\linewidth}}{
        \cellcolor[rgb]{0.9,0.9,0.9}{
            \makecell[{{p{\linewidth}}}]{
                \texttt{\tiny{[GM$|$GM]}}
                \texttt{<think>} \\
\texttt{Okay, let's try to figure this out. I need to divide the items between me and the other player. My values are plier:2, chalk:3, nail:2, bucket:1. The other player's values are secret. The key is that our proposals must be complementary, so we can't both want the same items. My goal is to maximize the total score for both of us.} \\
\\ 
\texttt{First, I should consider which items I value the most. Chalk is highest at 3, then pliers and nail at 2 each, and bucket at 1. Since there's only one chalk, if I take it, the other player can't. But maybe the other player values chalk more than I do. However, since I don't know their values, I have to make an educated guess.} \\
\\ 
\texttt{The items are 2 pliers, 1 chalk, 1 nail, 1 bucket. Let's think about possible splits. If I take the chalk (3) and one plier (2), that's 5 points for me. Then the other player could take the remaining plier, nail, and bucket. Their total would depend on their values for those items. If they value the plier, nail, and bucket highly, that could work. But if they don't, maybe they'd prefer something else.} \\
\\ 
\texttt{Alternatively, if I take the nail (2) and bucket (1), that's 3 points, but maybe the other player can take the chalk and pliers. But that might not be optimal for me. Wait, the total combined value would be higher if I take the chalk and plier, and they take the other plier, nail, and bucket. Let's see: my 3+2=5, and their plier (2) + nail (2) + bucket (1) = 5. Total 10. If I took nail and bucket (3) and they took chalk and pliers (3+2=5), total is 8. So the first option is better for total.} \\
\\ 
\texttt{But since I don't know their values, maybe they value the nail or bucket more. But without knowing, I have to assume that the other player's values are such that taking the remaining items after I take chalk and one plier would give them a good score. So maybe I should propose taking chalk and one plier. Then the other player can take the other} \\
            }
        }
    }
    & & \\ \\

    \theutterance \stepcounter{utterance}  
    & & \multicolumn{4}{p{0.6\linewidth}}{
        \cellcolor[rgb]{0.9,0.9,0.9}{
            \makecell[{{p{\linewidth}}}]{
                \texttt{\tiny{[GM$\rangle$P2]}}
                \texttt{<think>} \\
\texttt{Okay, let's try to figure this out. I need to divide the items between me and the other player. My values are plier:2, chalk:3, nail:2, bucket:1. The other player's values are secret. The key is that our proposals must be complementary, so we can't both want the same items. My goal is to maximize the total score for both of us.} \\
\\ 
\texttt{First, I should consider which items I value the most. Chalk is highest at 3, then pliers and nail at 2 each, and bucket at 1. Since there's only one chalk, if I take it, the other player can't. But maybe the other player values chalk more than I do. However, since I don't know their values, I have to make an educated guess.} \\
\\ 
\texttt{The items are 2 pliers, 1 chalk, 1 nail, 1 bucket. Let's think about possible splits. If I take the chalk (3) and one plier (2), that's 5 points for me. Then the other player could take the remaining plier, nail, and bucket. Their total would depend on their values for those items. If they value the plier, nail, and bucket highly, that could work. But if they don't, maybe they'd prefer something else.} \\
\\ 
\texttt{Alternatively, if I take the nail (2) and bucket (1), that's 3 points, but maybe the other player can take the chalk and pliers. But that might not be optimal for me. Wait, the total combined value would be higher if I take the chalk and plier, and they take the other plier, nail, and bucket. Let's see: my 3+2=5, and their plier (2) + nail (2) + bucket (1) = 5. Total 10. If I took nail and bucket (3) and they took chalk and pliers (3+2=5), total is 8. So the first option is better for total.} \\
\\ 
\texttt{But since I don't know their values, maybe they value the nail or bucket more. But without knowing, I have to assume that the other player's values are such that taking the remaining items after I take chalk and one plier would give them a good score. So maybe I should propose taking chalk and one plier. Then the other player can take the other} \\
            }
        }
    }
    & \\ \\

    \theutterance \stepcounter{utterance}  
    & & & \multicolumn{4}{p{0.6\linewidth}}{
        \cellcolor[rgb]{0.9,0.9,0.9}{
            \makecell[{{p{\linewidth}}}]{
                \texttt{\tiny{[GM$\langle$P2]}}
                \texttt{<think>} \\
\texttt{Okay, let's try to figure out the best way to split the items. My values are plier:2, chalk:3, nail:2, bucket:1. The other player's values are secret, so I need to make a proposal that allows for a complementary split.} \\
\\ 
\texttt{First, I should consider which items I value the most. Chalk is the highest at 3, then pliers and nail at 2 each, and bucket at 1. Since there's only one chalk, if I take it, the other player can't. But maybe the other player values chalk more than I do. However, since I don't know their values, I have to make an educated guess.} \\
\\ 
\texttt{The items are 2 pliers, 1 chalk, 1 nail, 1 bucket. Let's think about possible splits. If I take the chalk (3) and one plier (2), that's 5 points for me. Then the other player could take the remaining plier, nail, and bucket. Their total would depend on their values for those items. If they value the plier, nail, and bucket highly, that could work. But if they don't, maybe they'd prefer something else.} \\
\\ 
\texttt{Alternatively, if I take the nail (2) and bucket (1), that's 3 points, but maybe the other player can take the chalk and pliers. But that might not be optimal for me. Wait, the total combined value would be higher if I take the chalk and plier, and they take the other plier, nail, and bucket. Let's see: my 3+2=5, and their plier (2) + nail (2) + bucket (1) = 5. Total 10. If I took nail and bucket (3) and they took chalk and pliers (3+2=5), total is 8. So the first option is better for total.} \\
\\ 
\texttt{But since I don't know their values, maybe they value the nail or bucket more. But without knowing, I have to assume that the other player's values are such that taking the remaining items after I take chalk and one plier would give them a good score. So maybe I should propose taking chalk and one plier. Then the other player can take the other plier, nail, and bucket. That way, our proposals are complementary, and the total score is maximized.} \\
\\ 
\texttt{Alternatively} \\
            }
        }
    }
     \\ \\

    \theutterance \stepcounter{utterance}  
    & & & \multicolumn{2}{p{0.3\linewidth}}{
        \cellcolor[rgb]{0.9,0.9,0.9}{
            \makecell[{{p{\linewidth}}}]{
                \texttt{\tiny{[GM$|$GM]}}
                \texttt{continue} \\
            }
        }
    }
    & & \\ \\

    \theutterance \stepcounter{utterance}  
    & & & \multicolumn{2}{p{0.3\linewidth}}{
        \cellcolor[rgb]{0.9,0.9,0.9}{
            \makecell[{{p{\linewidth}}}]{
                \texttt{\tiny{[GM$|$GM]}}
                \texttt{<think>} \\
\texttt{Okay, let's try to figure out the best way to split the items. My values are plier:2, chalk:3, nail:2, bucket:1. The other player's values are secret, so I need to make a proposal that allows for a complementary split.} \\
\\ 
\texttt{First, I should consider which items I value the most. Chalk is the highest at 3, then pliers and nail at 2 each, and bucket at 1. Since there's only one chalk, if I take it, the other player can't. But maybe the other player values chalk more than I do. However, since I don't know their values, I have to make an educated guess.} \\
\\ 
\texttt{The items are 2 pliers, 1 chalk, 1 nail, 1 bucket. Let's think about possible splits. If I take the chalk (3) and one plier (2), that's 5 points for me. Then the other player could take the remaining plier, nail, and bucket. Their total would depend on their values for those items. If they value the plier, nail, and bucket highly, that could work. But if they don't, maybe they'd prefer something else.} \\
\\ 
\texttt{Alternatively, if I take the nail (2) and bucket (1), that's 3 points, but maybe the other player can take the chalk and pliers. But that might not be optimal for me. Wait, the total combined value would be higher if I take the chalk and plier, and they take the other plier, nail, and bucket. Let's see: my 3+2=5, and their plier (2) + nail (2) + bucket (1) = 5. Total 10. If I took nail and bucket (3) and they took chalk and pliers (3+2=5), total is 8. So the first option is better for total.} \\
\\ 
\texttt{But since I don't know their values, maybe they value the nail or bucket more. But without knowing, I have to assume that the other player's values are such that taking the remaining items after I take chalk and one plier would give them a good score. So maybe I should propose taking chalk and one plier. Then the other player can take the other plier, nail, and bucket. That way, our proposals are complementary, and the total score is maximized.} \\
\\ 
\texttt{Alternatively} \\
            }
        }
    }
    & & \\ \\

    \theutterance \stepcounter{utterance}  
    & & \multicolumn{4}{p{0.6\linewidth}}{
        \cellcolor[rgb]{0.9,0.9,0.9}{
            \makecell[{{p{\linewidth}}}]{
                \texttt{\tiny{[P1$\langle$GM]}}
                \texttt{<think>} \\
\texttt{Okay, let's try to figure out the best way to split the items. My values are plier:2, chalk:3, nail:2, bucket:1. The other player's values are secret, so I need to make a proposal that allows for a complementary split.} \\
\\ 
\texttt{First, I should consider which items I value the most. Chalk is the highest at 3, then pliers and nail at 2 each, and bucket at 1. Since there's only one chalk, if I take it, the other player can't. But maybe the other player values chalk more than I do. However, since I don't know their values, I have to make an educated guess.} \\
\\ 
\texttt{The items are 2 pliers, 1 chalk, 1 nail, 1 bucket. Let's think about possible splits. If I take the chalk (3) and one plier (2), that's 5 points for me. Then the other player could take the remaining plier, nail, and bucket. Their total would depend on their values for those items. If they value the plier, nail, and bucket highly, that could work. But if they don't, maybe they'd prefer something else.} \\
\\ 
\texttt{Alternatively, if I take the nail (2) and bucket (1), that's 3 points, but maybe the other player can take the chalk and pliers. But that might not be optimal for me. Wait, the total combined value would be higher if I take the chalk and plier, and they take the other plier, nail, and bucket. Let's see: my 3+2=5, and their plier (2) + nail (2) + bucket (1) = 5. Total 10. If I took nail and bucket (3) and they took chalk and pliers (3+2=5), total is 8. So the first option is better for total.} \\
\\ 
\texttt{But since I don't know their values, maybe they value the nail or bucket more. But without knowing, I have to assume that the other player's values are such that taking the remaining items after I take chalk and one plier would give them a good score. So maybe I should propose taking chalk and one plier. Then the other player can take the other plier, nail, and bucket. That way, our proposals are complementary, and the total score is maximized.} \\
\\ 
\texttt{Alternatively} \\
            }
        }
    }
    & \\ \\

    \theutterance \stepcounter{utterance}  
    & \multicolumn{4}{p{0.6\linewidth}}{
        \cellcolor[rgb]{0.9,0.9,0.9}{
            \makecell[{{p{\linewidth}}}]{
                \texttt{\tiny{[P1$\rangle$GM]}}
                \texttt{<think>} \\
\texttt{Okay, let's try to figure out the best way to split the items. My values are plier:2, chalk:3, nail:2, bucket:1. The other player's values are secret, so I need to make a proposal that allows for a complementary split.} \\
\\ 
\texttt{First, I should consider which items I value the most. Chalk is the highest at 3, then pliers and nail at 2 each, and bucket at 1. Since there's only one chalk, if I take it, the other player can't. But maybe the other player values chalk more than I do. However, since I don't know their values, I have to make an educated guess.} \\
\\ 
\texttt{The items are 2 pliers, 1 chalk, 1 nail, 1 bucket. Let's think about possible splits. If I take the chalk (3) and one plier (2), that's 5 points for me. Then the other player could take the remaining plier, nail, and bucket. Their total would depend on their values for those items. If they value the plier, nail, and bucket highly, that could work. But if they don't, maybe they'd prefer something else.} \\
\\ 
\texttt{Alternatively, if I take the nail (2) and bucket (1), that's 3 points, but maybe the other player can take the chalk and pliers. But that might not be optimal for me. Wait, the total combined value would be higher if I take the chalk and plier, and they take the other plier, nail, and bucket. Let's see: my 3+2=5, and their plier (2) + nail (2) + bucket (1) = 5. Total 10. If I took nail and bucket (3) and they took chalk and pliers (3+2=5), total is 8. So the first option is better for total.} \\
\\ 
\texttt{But since I don't know their values, maybe they value the nail or bucket more. But without knowing, I have to assume that the other player's values are such that taking the remaining items after I take chalk and one plier would give them a good score. So maybe I should propose taking chalk and one plier. Then the other player can take the other plier, nail, and bucket. That way, our proposals are complementary, and the total score is maximized.} \\
\\ 
\texttt{Alternatively} \\
            }
        }
    }
    & & \\ \\

    \theutterance \stepcounter{utterance}  
    & & & \multicolumn{2}{p{0.3\linewidth}}{
        \cellcolor[rgb]{0.9,0.9,0.9}{
            \makecell[{{p{\linewidth}}}]{
                \texttt{\tiny{[GM$|$GM]}}
                \texttt{continue} \\
            }
        }
    }
    & & \\ \\

    \theutterance \stepcounter{utterance}  
    & & & \multicolumn{2}{p{0.3\linewidth}}{
        \cellcolor[rgb]{0.9,0.9,0.9}{
            \makecell[{{p{\linewidth}}}]{
                \texttt{\tiny{[GM$|$GM]}}
                \texttt{<think>} \\
\texttt{Okay, let's try to figure out the best way to split the items. My values are plier:2, chalk:3, nail:2, bucket:1. The other player's values are secret, so I need to make a proposal that allows for a complementary split.} \\
\\ 
\texttt{First, I should consider which items I value the most. Chalk is the highest at 3, then pliers and nail at 2 each, and bucket at 1. Since there's only one chalk, if I take it, the other player can't. But maybe the other player values chalk more than I do. However, since I don't know their values, I have to make an educated guess.} \\
\\ 
\texttt{The items are 2 pliers, 1 chalk, 1 nail, 1 bucket. Let's think about possible splits. If I take the chalk (3) and one plier (2), that's 5 points for me. Then the other player could take the remaining plier, nail, and bucket. Their total would depend on their values for those items. If they value the plier, nail, and bucket highly, that could work. But if they don't, maybe they'd prefer something else.} \\
\\ 
\texttt{Alternatively, if I take the nail (2) and bucket (1), that's 3 points, but maybe the other player can take the chalk and pliers. But that might not be optimal for me. Wait, the total combined value would be higher if I take the chalk and plier, and they take the other plier, nail, and bucket. Let's see: my 3+2=5, and their plier (2) + nail (2) + bucket (1) = 5. Total 10. If I took nail and bucket (3) and they took chalk and pliers (3+2=5), total is 8. So the first option is better for total.} \\
\\ 
\texttt{But since I don't know their values, maybe they value the nail or bucket more. But without knowing, I have to assume that the other player's values are such that taking the remaining items after I take chalk and one plier would give them a good score. So maybe I should propose taking chalk and one plier. Then the other player can take the other plier, nail, and bucket. That way, our proposals are complementary, and the total score is maximized.} \\
\\ 
\texttt{Alternatively} \\
            }
        }
    }
    & & \\ \\

    \theutterance \stepcounter{utterance}  
    & & \multicolumn{4}{p{0.6\linewidth}}{
        \cellcolor[rgb]{0.9,0.9,0.9}{
            \makecell[{{p{\linewidth}}}]{
                \texttt{\tiny{[GM$\rangle$P2]}}
                \texttt{<think>} \\
\texttt{Okay, let's try to figure out the best way to split the items. My values are plier:2, chalk:3, nail:2, bucket:1. The other player's values are secret, so I need to make a proposal that allows for a complementary split.} \\
\\ 
\texttt{First, I should consider which items I value the most. Chalk is the highest at 3, then pliers and nail at 2 each, and bucket at 1. Since there's only one chalk, if I take it, the other player can't. But maybe the other player values chalk more than I do. However, since I don't know their values, I have to make an educated guess.} \\
\\ 
\texttt{The items are 2 pliers, 1 chalk, 1 nail, 1 bucket. Let's think about possible splits. If I take the chalk (3) and one plier (2), that's 5 points for me. Then the other player could take the remaining plier, nail, and bucket. Their total would depend on their values for those items. If they value the plier, nail, and bucket highly, that could work. But if they don't, maybe they'd prefer something else.} \\
\\ 
\texttt{Alternatively, if I take the nail (2) and bucket (1), that's 3 points, but maybe the other player can take the chalk and pliers. But that might not be optimal for me. Wait, the total combined value would be higher if I take the chalk and plier, and they take the other plier, nail, and bucket. Let's see: my 3+2=5, and their plier (2) + nail (2) + bucket (1) = 5. Total 10. If I took nail and bucket (3) and they took chalk and pliers (3+2=5), total is 8. So the first option is better for total.} \\
\\ 
\texttt{But since I don't know their values, maybe they value the nail or bucket more. But without knowing, I have to assume that the other player's values are such that taking the remaining items after I take chalk and one plier would give them a good score. So maybe I should propose taking chalk and one plier. Then the other player can take the other plier, nail, and bucket. That way, our proposals are complementary, and the total score is maximized.} \\
\\ 
\texttt{Alternatively} \\
            }
        }
    }
    & \\ \\

    \theutterance \stepcounter{utterance}  
    & & & \multicolumn{4}{p{0.6\linewidth}}{
        \cellcolor[rgb]{0.9,0.9,0.9}{
            \makecell[{{p{\linewidth}}}]{
                \texttt{\tiny{[GM$\langle$P2]}}
                \texttt{<think>} \\
\texttt{Okay, let's try to figure out the best way to split the items. My values are plier:2, chalk:3, nail:2, bucket:1. The other player's values are secret, so I need to make a proposal that allows for a complementary split.} \\
\\ 
\texttt{First, I should consider which items I value the most. Chalk is the highest at 3, then pliers and nail at 2 each, and bucket at 1. Since there's only one chalk, if I take it, the other player can't. But maybe the other player values chalk more than I do. However, since I don't know their values, I have to make an educated guess.} \\
\\ 
\texttt{The items are 2 pliers, 1 chalk, 1 nail, 1 bucket. Let's think about possible splits. If I take the chalk (3) and one plier (2), that's 5 points for me. Then the other player could take the remaining plier, nail, and bucket. Their total would depend on their values for those items. If they value the plier, nail, and bucket highly, that could work. But if they don't, maybe they'd prefer something else.} \\
\\ 
\texttt{Alternatively, if I take the nail (2) and bucket (1), that's 3 points, but maybe the other player can take the chalk and pliers. But that might not be optimal for me. Wait, the total combined value would be higher if I take the chalk and plier, and they take the other plier, nail, and bucket. Let's see: my 3+2=5, and their plier (2) + nail (2) + bucket (1) = 5. Total 10. If I took nail and bucket (3) and they took chalk and pliers (3+2=5), total is 8. So the first option is better for total.} \\
\\ 
\texttt{But since I don't know their values, maybe they value the nail or bucket more. But without knowing, I have to assume that the other player's values are such that taking the remaining items after I take chalk and one plier would give them a good score. So maybe I should propose taking chalk and one plier. Then the other player can take the other plier, nail, and bucket. That way, our proposals are complementary, and the total score is maximized.} \\
\\ 
\texttt{Alternatively} \\
            }
        }
    }
     \\ \\

    \theutterance \stepcounter{utterance}  
    & & & \multicolumn{2}{p{0.3\linewidth}}{
        \cellcolor[rgb]{0.9,0.9,0.9}{
            \makecell[{{p{\linewidth}}}]{
                \texttt{\tiny{[GM$|$GM]}}
                \texttt{continue} \\
            }
        }
    }
    & & \\ \\

    \theutterance \stepcounter{utterance}  
    & & & \multicolumn{2}{p{0.3\linewidth}}{
        \cellcolor[rgb]{0.9,0.9,0.9}{
            \makecell[{{p{\linewidth}}}]{
                \texttt{\tiny{[GM$|$GM]}}
                \texttt{<think>} \\
\texttt{Okay, let's try to figure out the best way to split the items. My values are plier:2, chalk:3, nail:2, bucket:1. The other player's values are secret, so I need to make a proposal that allows for a complementary split.} \\
\\ 
\texttt{First, I should consider which items I value the most. Chalk is the highest at 3, then pliers and nail at 2 each, and bucket at 1. Since there's only one chalk, if I take it, the other player can't. But maybe the other player values chalk more than I do. However, since I don't know their values, I have to make an educated guess.} \\
\\ 
\texttt{The items are 2 pliers, 1 chalk, 1 nail, 1 bucket. Let's think about possible splits. If I take the chalk (3) and one plier (2), that's 5 points for me. Then the other player could take the remaining plier, nail, and bucket. Their total would depend on their values for those items. If they value the plier, nail, and bucket highly, that could work. But if they don't, maybe they'd prefer something else.} \\
\\ 
\texttt{Alternatively, if I take the nail (2) and bucket (1), that's 3 points, but maybe the other player can take the chalk and pliers. But that might not be optimal for me. Wait, the total combined value would be higher if I take the chalk and plier, and they take the other plier, nail, and bucket. Let's see: my 3+2=5, and their plier (2) + nail (2) + bucket (1) = 5. Total 10. If I took nail and bucket (3) and they took chalk and pliers (3+2=5), total is 8. So the first option is better for total.} \\
\\ 
\texttt{But since I don't know their values, maybe they value the nail or bucket more. But without knowing, I have to assume that the other player's values are such that taking the remaining items after I take chalk and one plier would give them a good score. So maybe I should propose taking chalk and one plier. Then the other player can take the other plier, nail, and bucket. That way, our proposals are complementary, and the total score is maximized.} \\
\\ 
\texttt{Alternatively} \\
            }
        }
    }
    & & \\ \\

    \theutterance \stepcounter{utterance}  
    & & \multicolumn{4}{p{0.6\linewidth}}{
        \cellcolor[rgb]{0.9,0.9,0.9}{
            \makecell[{{p{\linewidth}}}]{
                \texttt{\tiny{[P1$\langle$GM]}}
                \texttt{<think>} \\
\texttt{Okay, let's try to figure out the best way to split the items. My values are plier:2, chalk:3, nail:2, bucket:1. The other player's values are secret, so I need to make a proposal that allows for a complementary split.} \\
\\ 
\texttt{First, I should consider which items I value the most. Chalk is the highest at 3, then pliers and nail at 2 each, and bucket at 1. Since there's only one chalk, if I take it, the other player can't. But maybe the other player values chalk more than I do. However, since I don't know their values, I have to make an educated guess.} \\
\\ 
\texttt{The items are 2 pliers, 1 chalk, 1 nail, 1 bucket. Let's think about possible splits. If I take the chalk (3) and one plier (2), that's 5 points for me. Then the other player could take the remaining plier, nail, and bucket. Their total would depend on their values for those items. If they value the plier, nail, and bucket highly, that could work. But if they don't, maybe they'd prefer something else.} \\
\\ 
\texttt{Alternatively, if I take the nail (2) and bucket (1), that's 3 points, but maybe the other player can take the chalk and pliers. But that might not be optimal for me. Wait, the total combined value would be higher if I take the chalk and plier, and they take the other plier, nail, and bucket. Let's see: my 3+2=5, and their plier (2) + nail (2) + bucket (1) = 5. Total 10. If I took nail and bucket (3) and they took chalk and pliers (3+2=5), total is 8. So the first option is better for total.} \\
\\ 
\texttt{But since I don't know their values, maybe they value the nail or bucket more. But without knowing, I have to assume that the other player's values are such that taking the remaining items after I take chalk and one plier would give them a good score. So maybe I should propose taking chalk and one plier. Then the other player can take the other plier, nail, and bucket. That way, our proposals are complementary, and the total score is maximized.} \\
\\ 
\texttt{Alternatively} \\
            }
        }
    }
    & \\ \\

    \theutterance \stepcounter{utterance}  
    & \multicolumn{4}{p{0.6\linewidth}}{
        \cellcolor[rgb]{0.9,0.9,0.9}{
            \makecell[{{p{\linewidth}}}]{
                \texttt{\tiny{[P1$\rangle$GM]}}
                \texttt{<think>} \\
\texttt{Okay, let's try to figure out the best way to split the items. My values are plier:2, chalk:3, nail:2, bucket:1. The other player's values are secret, so I need to make a proposal that allows for a complementary split.} \\
\\ 
\texttt{First, I should consider which items I value the most. Chalk is the highest at 3, then pliers and nail at 2 each, and bucket at 1. Since there's only one chalk, if I take it, the other player can't. But maybe the other player values chalk more than I do. However, since I don't know their values, I have to make an educated guess.} \\
\\ 
\texttt{The items are 2 pliers, 1 chalk, 1 nail, 1 bucket. Let's think about possible splits. If I take the chalk (3) and one plier (2), that's 5 points for me. Then the other player could take the remaining plier, nail, and bucket. Their total would depend on their values for those items. If they value the plier, nail, and bucket highly, that could work. But if they don't, maybe they'd prefer something else.} \\
\\ 
\texttt{Alternatively, if I take the nail (2) and bucket (1), that's 3 points, but maybe the other player can take the chalk and pliers. But that might not be optimal for me. Wait, the total combined value would be higher if I take the chalk and plier, and they take the other plier, nail, and bucket. Let's see: my 3+2=5, and their plier (2) + nail (2) + bucket (1) = 5. Total 10. If I took nail and bucket (3) and they took chalk and pliers (3+2=5), total is 8. So the first option is better for total.} \\
\\ 
\texttt{But since I don't know their values, maybe they value the nail or bucket more. But without knowing, I have to assume that the other player's values are such that taking the remaining items after I take chalk and one plier would give them a good score. So maybe I should propose taking chalk and one plier. Then the other player can take the other plier, nail, and bucket. That way, our proposals are complementary, and the total score is maximized.} \\
\\ 
\texttt{Alternatively} \\
            }
        }
    }
    & & \\ \\

    \theutterance \stepcounter{utterance}  
    & & & \multicolumn{2}{p{0.3\linewidth}}{
        \cellcolor[rgb]{0.9,0.9,0.9}{
            \makecell[{{p{\linewidth}}}]{
                \texttt{\tiny{[GM$|$GM]}}
                \texttt{continue} \\
            }
        }
    }
    & & \\ \\

    \theutterance \stepcounter{utterance}  
    & & & \multicolumn{2}{p{0.3\linewidth}}{
        \cellcolor[rgb]{0.9,0.9,0.9}{
            \makecell[{{p{\linewidth}}}]{
                \texttt{\tiny{[GM$|$GM]}}
                \texttt{<think>} \\
\texttt{Okay, let's try to figure out the best way to split the items. My values are plier:2, chalk:3, nail:2, bucket:1. The other player's values are secret, so I need to make a proposal that allows for a complementary split.} \\
\\ 
\texttt{First, I should consider which items I value the most. Chalk is the highest at 3, then pliers and nail at 2 each, and bucket at 1. Since there's only one chalk, if I take it, the other player can't. But maybe the other player values chalk more than I do. However, since I don't know their values, I have to make an educated guess.} \\
\\ 
\texttt{The items are 2 pliers, 1 chalk, 1 nail, 1 bucket. Let's think about possible splits. If I take the chalk (3) and one plier (2), that's 5 points for me. Then the other player could take the remaining plier, nail, and bucket. Their total would depend on their values for those items. If they value the plier, nail, and bucket highly, that could work. But if they don't, maybe they'd prefer something else.} \\
\\ 
\texttt{Alternatively, if I take the nail (2) and bucket (1), that's 3 points, but maybe the other player can take the chalk and pliers. But that might not be optimal for me. Wait, the total combined value would be higher if I take the chalk and plier, and they take the other plier, nail, and bucket. Let's see: my 3+2=5, and their plier (2) + nail (2) + bucket (1) = 5. Total 10. If I took nail and bucket (3) and they took chalk and pliers (3+2=5), total is 8. So the first option is better for total.} \\
\\ 
\texttt{But since I don't know their values, maybe they value the nail or bucket more. But without knowing, I have to assume that the other player's values are such that taking the remaining items after I take chalk and one plier would give them a good score. So maybe I should propose taking chalk and one plier. Then the other player can take the other plier, nail, and bucket. That way, our proposals are complementary, and the total score is maximized.} \\
\\ 
\texttt{Alternatively} \\
            }
        }
    }
    & & \\ \\

    \theutterance \stepcounter{utterance}  
    & & \multicolumn{4}{p{0.6\linewidth}}{
        \cellcolor[rgb]{0.9,0.9,0.9}{
            \makecell[{{p{\linewidth}}}]{
                \texttt{\tiny{[GM$\rangle$P2]}}
                \texttt{<think>} \\
\texttt{Okay, let's try to figure out the best way to split the items. My values are plier:2, chalk:3, nail:2, bucket:1. The other player's values are secret, so I need to make a proposal that allows for a complementary split.} \\
\\ 
\texttt{First, I should consider which items I value the most. Chalk is the highest at 3, then pliers and nail at 2 each, and bucket at 1. Since there's only one chalk, if I take it, the other player can't. But maybe the other player values chalk more than I do. However, since I don't know their values, I have to make an educated guess.} \\
\\ 
\texttt{The items are 2 pliers, 1 chalk, 1 nail, 1 bucket. Let's think about possible splits. If I take the chalk (3) and one plier (2), that's 5 points for me. Then the other player could take the remaining plier, nail, and bucket. Their total would depend on their values for those items. If they value the plier, nail, and bucket highly, that could work. But if they don't, maybe they'd prefer something else.} \\
\\ 
\texttt{Alternatively, if I take the nail (2) and bucket (1), that's 3 points, but maybe the other player can take the chalk and pliers. But that might not be optimal for me. Wait, the total combined value would be higher if I take the chalk and plier, and they take the other plier, nail, and bucket. Let's see: my 3+2=5, and their plier (2) + nail (2) + bucket (1) = 5. Total 10. If I took nail and bucket (3) and they took chalk and pliers (3+2=5), total is 8. So the first option is better for total.} \\
\\ 
\texttt{But since I don't know their values, maybe they value the nail or bucket more. But without knowing, I have to assume that the other player's values are such that taking the remaining items after I take chalk and one plier would give them a good score. So maybe I should propose taking chalk and one plier. Then the other player can take the other plier, nail, and bucket. That way, our proposals are complementary, and the total score is maximized.} \\
\\ 
\texttt{Alternatively} \\
            }
        }
    }
    & \\ \\

    \theutterance \stepcounter{utterance}  
    & & & \multicolumn{4}{p{0.6\linewidth}}{
        \cellcolor[rgb]{0.9,0.9,0.9}{
            \makecell[{{p{\linewidth}}}]{
                \texttt{\tiny{[GM$\langle$P2]}}
                \texttt{<think>} \\
\texttt{Okay, let's try to figure out the best way to split the items. My values are plier:2, chalk:3, nail:2, bucket:1. The other player's values are secret, so I need to make a proposal that allows for a complementary split.} \\
\\ 
\texttt{First, I should consider which items I value the most. Chalk is the highest at 3, then pliers and nail at 2 each, and bucket at 1. Since there's only one chalk, if I take it, the other player can't. But maybe the other player values chalk more than I do. However, since I don't know their values, I have to make an educated guess.} \\
\\ 
\texttt{The items are 2 pliers, 1 chalk, 1 nail, 1 bucket. Let's think about possible splits. If I take the chalk (3) and one plier (2), that's 5 points for me. Then the other player could take the remaining plier, nail, and bucket. Their total would depend on their values for those items. If they value the plier, nail, and bucket highly, that could work. But if they don't, maybe they'd prefer something else.} \\
\\ 
\texttt{Alternatively, if I take the nail (2) and bucket (1), that's 3 points, but maybe the other player can take the chalk and pliers. But that might not be optimal for me. Wait, the total combined value would be higher if I take the chalk and plier, and they take the other plier, nail, and bucket. Let's see: my 3+2=5, and their plier (2) + nail (2) + bucket (1) = 5. Total 10. If I took nail and bucket (3) and they took chalk and pliers (3+2=5), total is 8. So the first option is better for total.} \\
\\ 
\texttt{But since I don't know their values, maybe they value the nail or bucket more. But without knowing, I have to assume that the other player's values are such that taking the remaining items after I take chalk and one plier would give them a good score. So maybe I should propose taking chalk and one plier. Then the other player can take the other plier, nail, and bucket. That way, our proposals are complementary, and the total score is maximized.} \\
\\ 
\texttt{Alternatively} \\
            }
        }
    }
     \\ \\

    \theutterance \stepcounter{utterance}  
    & & & \multicolumn{2}{p{0.3\linewidth}}{
        \cellcolor[rgb]{0.9,0.9,0.9}{
            \makecell[{{p{\linewidth}}}]{
                \texttt{\tiny{[GM$|$GM]}}
                \texttt{continue} \\
            }
        }
    }
    & & \\ \\

    \theutterance \stepcounter{utterance}  
    & & & \multicolumn{2}{p{0.3\linewidth}}{
        \cellcolor[rgb]{0.9,0.9,0.9}{
            \makecell[{{p{\linewidth}}}]{
                \texttt{\tiny{[GM$|$GM]}}
                \texttt{<think>} \\
\texttt{Okay, let's try to figure out the best way to split the items. My values are plier:2, chalk:3, nail:2, bucket:1. The other player's values are secret, so I need to make a proposal that allows for a complementary split.} \\
\\ 
\texttt{First, I should consider which items I value the most. Chalk is the highest at 3, then pliers and nail at 2 each, and bucket at 1. Since there's only one chalk, if I take it, the other player can't. But maybe the other player values chalk more than I do. However, since I don't know their values, I have to make an educated guess.} \\
\\ 
\texttt{The items are 2 pliers, 1 chalk, 1 nail, 1 bucket. Let's think about possible splits. If I take the chalk (3) and one plier (2), that's 5 points for me. Then the other player could take the remaining plier, nail, and bucket. Their total would depend on their values for those items. If they value the plier, nail, and bucket highly, that could work. But if they don't, maybe they'd prefer something else.} \\
\\ 
\texttt{Alternatively, if I take the nail (2) and bucket (1), that's 3 points, but maybe the other player can take the chalk and pliers. But that might not be optimal for me. Wait, the total combined value would be higher if I take the chalk and plier, and they take the other plier, nail, and bucket. Let's see: my 3+2=5, and their plier (2) + nail (2) + bucket (1) = 5. Total 10. If I took nail and bucket (3) and they took chalk and pliers (3+2=5), total is 8. So the first option is better for total.} \\
\\ 
\texttt{But since I don't know their values, maybe they value the nail or bucket more. But without knowing, I have to assume that the other player's values are such that taking the remaining items after I take chalk and one plier would give them a good score. So maybe I should propose taking chalk and one plier. Then the other player can take the other plier, nail, and bucket. That way, our proposals are complementary, and the total score is maximized.} \\
\\ 
\texttt{Alternatively} \\
            }
        }
    }
    & & \\ \\

    \theutterance \stepcounter{utterance}  
    & & \multicolumn{4}{p{0.6\linewidth}}{
        \cellcolor[rgb]{0.9,0.9,0.9}{
            \makecell[{{p{\linewidth}}}]{
                \texttt{\tiny{[P1$\langle$GM]}}
                \texttt{<think>} \\
\texttt{Okay, let's try to figure out the best way to split the items. My values are plier:2, chalk:3, nail:2, bucket:1. The other player's values are secret, so I need to make a proposal that allows for a complementary split.} \\
\\ 
\texttt{First, I should consider which items I value the most. Chalk is the highest at 3, then pliers and nail at 2 each, and bucket at 1. Since there's only one chalk, if I take it, the other player can't. But maybe the other player values chalk more than I do. However, since I don't know their values, I have to make an educated guess.} \\
\\ 
\texttt{The items are 2 pliers, 1 chalk, 1 nail, 1 bucket. Let's think about possible splits. If I take the chalk (3) and one plier (2), that's 5 points for me. Then the other player could take the remaining plier, nail, and bucket. Their total would depend on their values for those items. If they value the plier, nail, and bucket highly, that could work. But if they don't, maybe they'd prefer something else.} \\
\\ 
\texttt{Alternatively, if I take the nail (2) and bucket (1), that's 3 points, but maybe the other player can take the chalk and pliers. But that might not be optimal for me. Wait, the total combined value would be higher if I take the chalk and plier, and they take the other plier, nail, and bucket. Let's see: my 3+2=5, and their plier (2) + nail (2) + bucket (1) = 5. Total 10. If I took nail and bucket (3) and they took chalk and pliers (3+2=5), total is 8. So the first option is better for total.} \\
\\ 
\texttt{But since I don't know their values, maybe they value the nail or bucket more. But without knowing, I have to assume that the other player's values are such that taking the remaining items after I take chalk and one plier would give them a good score. So maybe I should propose taking chalk and one plier. Then the other player can take the other plier, nail, and bucket. That way, our proposals are complementary, and the total score is maximized.} \\
\\ 
\texttt{Alternatively} \\
            }
        }
    }
    & \\ \\

    \theutterance \stepcounter{utterance}  
    & \multicolumn{4}{p{0.6\linewidth}}{
        \cellcolor[rgb]{0.9,0.9,0.9}{
            \makecell[{{p{\linewidth}}}]{
                \texttt{\tiny{[P1$\rangle$GM]}}
                \texttt{<think>} \\
\texttt{Okay, let's try to figure out the best way to split the items. My values are plier:2, chalk:3, nail:2, bucket:1. The other player's values are secret, so I need to make a proposal that allows for a complementary split.} \\
\\ 
\texttt{First, I should consider which items I value the most. Chalk is the highest at 3, then pliers and nail at 2 each, and bucket at 1. Since there's only one chalk, if I take it, the other player can't. But maybe the other player values chalk more than I do. However, since I don't know their values, I have to make an educated guess.} \\
\\ 
\texttt{The items are 2 pliers, 1 chalk, 1 nail, 1 bucket. Let's think about possible splits. If I take the chalk (3) and one plier (2), that's 5 points for me. Then the other player could take the remaining plier, nail, and bucket. Their total would depend on their values for those items. If they value the plier, nail, and bucket highly, that could work. But if they don't, maybe they'd prefer something else.} \\
\\ 
\texttt{Alternatively, if I take the nail (2) and bucket (1), that's 3 points, but maybe the other player can take the chalk and pliers. But that might not be optimal for me. Wait, the total combined value would be higher if I take the chalk and plier, and they take the other plier, nail, and bucket. Let's see: my 3+2=5, and their plier (2) + nail (2) + bucket (1) = 5. Total 10. If I took nail and bucket (3) and they took chalk and pliers (3+2=5), total is 8. So the first option is better for total.} \\
\\ 
\texttt{But since I don't know their values, maybe they value the nail or bucket more. But without knowing, I have to assume that the other player's values are such that taking the remaining items after I take chalk and one plier would give them a good score. So maybe I should propose taking chalk and one plier. Then the other player can take the other plier, nail, and bucket. That way, our proposals are complementary, and the total score is maximized.} \\
\\ 
\texttt{Alternatively} \\
            }
        }
    }
    & & \\ \\

    \theutterance \stepcounter{utterance}  
    & & & \multicolumn{2}{p{0.3\linewidth}}{
        \cellcolor[rgb]{0.9,0.9,0.9}{
            \makecell[{{p{\linewidth}}}]{
                \texttt{\tiny{[GM$|$GM]}}
                \texttt{continue} \\
            }
        }
    }
    & & \\ \\

    \theutterance \stepcounter{utterance}  
    & & & \multicolumn{2}{p{0.3\linewidth}}{
        \cellcolor[rgb]{0.9,0.9,0.9}{
            \makecell[{{p{\linewidth}}}]{
                \texttt{\tiny{[GM$|$GM]}}
                \texttt{<think>} \\
\texttt{Okay, let's try to figure out the best way to split the items. My values are plier:2, chalk:3, nail:2, bucket:1. The other player's values are secret, so I need to make a proposal that allows for a complementary split.} \\
\\ 
\texttt{First, I should consider which items I value the most. Chalk is the highest at 3, then pliers and nail at 2 each, and bucket at 1. Since there's only one chalk, if I take it, the other player can't. But maybe the other player values chalk more than I do. However, since I don't know their values, I have to make an educated guess.} \\
\\ 
\texttt{The items are 2 pliers, 1 chalk, 1 nail, 1 bucket. Let's think about possible splits. If I take the chalk (3) and one plier (2), that's 5 points for me. Then the other player could take the remaining plier, nail, and bucket. Their total would depend on their values for those items. If they value the plier, nail, and bucket highly, that could work. But if they don't, maybe they'd prefer something else.} \\
\\ 
\texttt{Alternatively, if I take the nail (2) and bucket (1), that's 3 points, but maybe the other player can take the chalk and pliers. But that might not be optimal for me. Wait, the total combined value would be higher if I take the chalk and plier, and they take the other plier, nail, and bucket. Let's see: my 3+2=5, and their plier (2) + nail (2) + bucket (1) = 5. Total 10. If I took nail and bucket (3) and they took chalk and pliers (3+2=5), total is 8. So the first option is better for total.} \\
\\ 
\texttt{But since I don't know their values, maybe they value the nail or bucket more. But without knowing, I have to assume that the other player's values are such that taking the remaining items after I take chalk and one plier would give them a good score. So maybe I should propose taking chalk and one plier. Then the other player can take the other plier, nail, and bucket. That way, our proposals are complementary, and the total score is maximized.} \\
\\ 
\texttt{Alternatively} \\
            }
        }
    }
    & & \\ \\

    \theutterance \stepcounter{utterance}  
    & & \multicolumn{4}{p{0.6\linewidth}}{
        \cellcolor[rgb]{0.9,0.9,0.9}{
            \makecell[{{p{\linewidth}}}]{
                \texttt{\tiny{[GM$\rangle$P2]}}
                \texttt{<think>} \\
\texttt{Okay, let's try to figure out the best way to split the items. My values are plier:2, chalk:3, nail:2, bucket:1. The other player's values are secret, so I need to make a proposal that allows for a complementary split.} \\
\\ 
\texttt{First, I should consider which items I value the most. Chalk is the highest at 3, then pliers and nail at 2 each, and bucket at 1. Since there's only one chalk, if I take it, the other player can't. But maybe the other player values chalk more than I do. However, since I don't know their values, I have to make an educated guess.} \\
\\ 
\texttt{The items are 2 pliers, 1 chalk, 1 nail, 1 bucket. Let's think about possible splits. If I take the chalk (3) and one plier (2), that's 5 points for me. Then the other player could take the remaining plier, nail, and bucket. Their total would depend on their values for those items. If they value the plier, nail, and bucket highly, that could work. But if they don't, maybe they'd prefer something else.} \\
\\ 
\texttt{Alternatively, if I take the nail (2) and bucket (1), that's 3 points, but maybe the other player can take the chalk and pliers. But that might not be optimal for me. Wait, the total combined value would be higher if I take the chalk and plier, and they take the other plier, nail, and bucket. Let's see: my 3+2=5, and their plier (2) + nail (2) + bucket (1) = 5. Total 10. If I took nail and bucket (3) and they took chalk and pliers (3+2=5), total is 8. So the first option is better for total.} \\
\\ 
\texttt{But since I don't know their values, maybe they value the nail or bucket more. But without knowing, I have to assume that the other player's values are such that taking the remaining items after I take chalk and one plier would give them a good score. So maybe I should propose taking chalk and one plier. Then the other player can take the other plier, nail, and bucket. That way, our proposals are complementary, and the total score is maximized.} \\
\\ 
\texttt{Alternatively} \\
            }
        }
    }
    & \\ \\

    \theutterance \stepcounter{utterance}  
    & & & \multicolumn{4}{p{0.6\linewidth}}{
        \cellcolor[rgb]{0.9,0.9,0.9}{
            \makecell[{{p{\linewidth}}}]{
                \texttt{\tiny{[GM$\langle$P2]}}
                \texttt{<think>} \\
\texttt{Okay, let's try to figure out the best way to split the items. My values are plier:2, chalk:3, nail:2, bucket:1. The other player's values are secret, so I need to make a proposal that allows for a complementary split.} \\
\\ 
\texttt{First, I should consider which items I value the most. Chalk is the highest at 3, then pliers and nail at 2 each, and bucket at 1. Since there's only one chalk, if I take it, the other player can't. But maybe the other player values chalk more than I do. However, since I don't know their values, I have to make an educated guess.} \\
\\ 
\texttt{The items are 2 pliers, 1 chalk, 1 nail, 1 bucket. Let's think about possible splits. If I take the chalk (3) and one plier (2), that's 5 points for me. Then the other player could take the remaining plier, nail, and bucket. Their total would depend on their values for those items. If they value the plier, nail, and bucket highly, that could work. But if they don't, maybe they'd prefer something else.} \\
\\ 
\texttt{Alternatively, if I take the nail (2) and bucket (1), that's 3 points, but maybe the other player can take the chalk and pliers. But that might not be optimal for me. Wait, the total combined value would be higher if I take the chalk and plier, and they take the other plier, nail, and bucket. Let's see: my 3+2=5, and their plier (2) + nail (2) + bucket (1) = 5. Total 10. If I took nail and bucket (3) and they took chalk and pliers (3+2=5), total is 8. So the first option is better for total.} \\
\\ 
\texttt{But since I don't know their values, maybe they value the nail or bucket more. But without knowing, I have to assume that the other player's values are such that taking the remaining items after I take chalk and one plier would give them a good score. So maybe I should propose taking chalk and one plier. Then the other player can take the other plier, nail, and bucket. That way, our proposals are complementary, and the total score is maximized.} \\
\\ 
\texttt{Alternatively} \\
            }
        }
    }
     \\ \\

    \theutterance \stepcounter{utterance}  
    & & & \multicolumn{2}{p{0.3\linewidth}}{
        \cellcolor[rgb]{0.9,0.9,0.9}{
            \makecell[{{p{\linewidth}}}]{
                \texttt{\tiny{[GM$|$GM]}}
                \texttt{continue} \\
            }
        }
    }
    & & \\ \\

    \theutterance \stepcounter{utterance}  
    & & & \multicolumn{2}{p{0.3\linewidth}}{
        \cellcolor[rgb]{0.9,0.9,0.9}{
            \makecell[{{p{\linewidth}}}]{
                \texttt{\tiny{[GM$|$GM]}}
                \texttt{<think>} \\
\texttt{Okay, let's try to figure out the best way to split the items. My values are plier:2, chalk:3, nail:2, bucket:1. The other player's values are secret, so I need to make a proposal that allows for a complementary split.} \\
\\ 
\texttt{First, I should consider which items I value the most. Chalk is the highest at 3, then pliers and nail at 2 each, and bucket at 1. Since there's only one chalk, if I take it, the other player can't. But maybe the other player values chalk more than I do. However, since I don't know their values, I have to make an educated guess.} \\
\\ 
\texttt{The items are 2 pliers, 1 chalk, 1 nail, 1 bucket. Let's think about possible splits. If I take the chalk (3) and one plier (2), that's 5 points for me. Then the other player could take the remaining plier, nail, and bucket. Their total would depend on their values for those items. If they value the plier, nail, and bucket highly, that could work. But if they don't, maybe they'd prefer something else.} \\
\\ 
\texttt{Alternatively, if I take the nail (2) and bucket (1), that's 3 points, but maybe the other player can take the chalk and pliers. But that might not be optimal for me. Wait, the total combined value would be higher if I take the chalk and plier, and they take the other plier, nail, and bucket. Let's see: my 3+2=5, and their plier (2) + nail (2) + bucket (1) = 5. Total 10. If I took nail and bucket (3) and they took chalk and pliers (3+2=5), total is 8. So the first option is better for total.} \\
\\ 
\texttt{But since I don't know their values, maybe they value the nail or bucket more. But without knowing, I have to assume that the other player's values are such that taking the remaining items after I take chalk and one plier would give them a good score. So maybe I should propose taking chalk and one plier. Then the other player can take the other plier, nail, and bucket. That way, our proposals are complementary, and the total score is maximized.} \\
\\ 
\texttt{Alternatively} \\
            }
        }
    }
    & & \\ \\

    \theutterance \stepcounter{utterance}  
    & & \multicolumn{4}{p{0.6\linewidth}}{
        \cellcolor[rgb]{0.9,0.9,0.9}{
            \makecell[{{p{\linewidth}}}]{
                \texttt{\tiny{[P1$\langle$GM]}}
                \texttt{<think>} \\
\texttt{Okay, let's try to figure out the best way to split the items. My values are plier:2, chalk:3, nail:2, bucket:1. The other player's values are secret, so I need to make a proposal that allows for a complementary split.} \\
\\ 
\texttt{First, I should consider which items I value the most. Chalk is the highest at 3, then pliers and nail at 2 each, and bucket at 1. Since there's only one chalk, if I take it, the other player can't. But maybe the other player values chalk more than I do. However, since I don't know their values, I have to make an educated guess.} \\
\\ 
\texttt{The items are 2 pliers, 1 chalk, 1 nail, 1 bucket. Let's think about possible splits. If I take the chalk (3) and one plier (2), that's 5 points for me. Then the other player could take the remaining plier, nail, and bucket. Their total would depend on their values for those items. If they value the plier, nail, and bucket highly, that could work. But if they don't, maybe they'd prefer something else.} \\
\\ 
\texttt{Alternatively, if I take the nail (2) and bucket (1), that's 3 points, but maybe the other player can take the chalk and pliers. But that might not be optimal for me. Wait, the total combined value would be higher if I take the chalk and plier, and they take the other plier, nail, and bucket. Let's see: my 3+2=5, and their plier (2) + nail (2) + bucket (1) = 5. Total 10. If I took nail and bucket (3) and they took chalk and pliers (3+2=5), total is 8. So the first option is better for total.} \\
\\ 
\texttt{But since I don't know their values, maybe they value the nail or bucket more. But without knowing, I have to assume that the other player's values are such that taking the remaining items after I take chalk and one plier would give them a good score. So maybe I should propose taking chalk and one plier. Then the other player can take the other plier, nail, and bucket. That way, our proposals are complementary, and the total score is maximized.} \\
\\ 
\texttt{Alternatively} \\
\\ 
\\ 
\texttt{The time is up. It is now your turn to submit a secret proposal.} \\
            }
        }
    }
    & \\ \\

    \theutterance \stepcounter{utterance}  
    & \multicolumn{4}{p{0.6\linewidth}}{
        \cellcolor[rgb]{0.9,0.9,0.9}{
            \makecell[{{p{\linewidth}}}]{
                \texttt{\tiny{[P1$\rangle$GM]}}
                \texttt{<think>} \\
\texttt{Okay, I need to make a secret proposal here. Let me recap the items and my values. The available items are 2 pliers, 1 chalk, 1 nail, and 1 bucket. My value function is plier:2, chalk:3, nail:2, bucket:1. The other player's values are unknown, but I need to propose a division that's complementary to theirs.} \\
\\ 
\texttt{Since my highest value is chalk at 3, I should prioritize getting that. Then, pliers and nail are next at 2 each. The bucket is the least valuable for me. If I take the chalk and one plier, that gives me 3 + 2 = 5 points. The remaining items would be one plier, nail, and bucket. If the other player takes those, their total would depend on their values, but since I don't know them, I have to assume they can get a decent score. The total combined value would be 5 (mine) plus whatever they get from the remaining items. } \\
\\ 
\texttt{Alternatively, if I take the nail and bucket, that's 2 + 1 = 3 points, but then the other player could take chalk and both pliers, which might give them a higher total. But since I don't know their values, maybe they value the chalk more. However, taking the chalk for myself seems better for my score. } \\
\\ 
\texttt{Another option is to take one plier and the nail, which is 2 + 2 = 4. Then the other player could take chalk, the other plier, and bucket. But again, without knowing their values, it's hard to say. } \\
\\ 
\texttt{The key is to make sure our proposals don't overlap. If I take chalk and one plier, the other player can take the remaining plier, nail, and bucket. That way, there's no overlap. The total combined value would be 5 (mine) plus their value for plier + nail + bucket. Since I don't know their values, I have to assume that this split allows for a complementary proposal. } \\
\\ 
\texttt{I think the best move is to propose taking the chalk and one plier. That way, I maximize my own score and leave the other items for the other player, which might align with their values. So my secret proposal would be {[}Proposal: 1 chalk, 1 plier{]}.} \\
\texttt{</think>} \\
\\ 
\texttt{{[}Proposal: 1 chalk,} \\
            }
        }
    }
    & & \\ \\

    \theutterance \stepcounter{utterance}  
    & & & \multicolumn{2}{p{0.3\linewidth}}{
        \cellcolor[rgb]{0.9,0.9,0.9}{
            \makecell[{{p{\linewidth}}}]{
                \texttt{\tiny{[GM$|$GM]}}
                \texttt{proposal} \\
            }
        }
    }
    & & \\ \\

    \theutterance \stepcounter{utterance}  
    & & & \multicolumn{2}{p{0.3\linewidth}}{
        \cellcolor[rgb]{0.9,0.9,0.9}{
            \makecell[{{p{\linewidth}}}]{
                \texttt{\tiny{[GM$|$GM]}}
                [1, 1, 0, 0]
            }
        }
    }
    & & \\ \\

    \theutterance \stepcounter{utterance}  
    & & \multicolumn{4}{p{0.6\linewidth}}{
        \cellcolor[rgb]{0.9,0.9,0.9}{
            \makecell[{{p{\linewidth}}}]{
                \texttt{\tiny{[GM$\rangle$P2]}}
                \texttt{The other player has submitted a secret proposal. It is now your turn to submit a proposal of your own.} \\
            }
        }
    }
    & \\ \\

    \theutterance \stepcounter{utterance}  
    & & & \multicolumn{4}{p{0.6\linewidth}}{
        \cellcolor[rgb]{0.9,0.9,0.9}{
            \makecell[{{p{\linewidth}}}]{
                \texttt{\tiny{[GM$\langle$P2]}}
                \texttt{<think>} \\
\texttt{Okay, the other player has already submitted a secret proposal, and now I need to submit mine. Let me think about what they might have proposed. Since their values are secret, I don't know their preferences, but I can try to deduce based on the items available.} \\
\\ 
\texttt{The available items are 2 pliers, 1 chalk, 1 nail, 1 bucket. My values are plier:2, chalk:3, nail:2, bucket:1. The other player's values are unknown. The goal is to have complementary proposals so that both can get their items without overlap.} \\
\\ 
\texttt{If the other player proposed taking some items, I need to make sure my proposal doesn't conflict. Since I don't know their proposal, I have to make a safe bet. The best way to maximize the total score is to split the items in a way that covers all high{-}value items for both players.} \\
\\ 
\texttt{Since I value chalk the most (3), I should try to get that. If I take chalk, then the other player can't take it. But maybe they also value it. However, without knowing, I have to assume they might not. So taking chalk and one plier (which I value at 2) gives me 5 points. Then the other player could take the remaining plier, nail, and bucket. If their values for those items are high, that works. If not, maybe they have a different split.} \\
\\ 
\texttt{Alternatively, if I take the nail and bucket, which are lower value for me, but maybe the other player wants the chalk and pliers. But that might not be optimal for me. The total combined value would be higher if I take the chalk and a plier, and they take the rest. Let's calculate: my 3+2=5, and if they take plier (2), nail (2), bucket (1), total is 5. Combined 10. If I took nail and bucket (3) and they took chalk and pliers (3+2=5), total is 8. So the first option is better.} \\
\\ 
\texttt{Therefore, I should propose taking chalk and one plier. That way, the other player can take the remaining items, and our proposals are complementary. Let me check if that's possible. If I take 1 chalk and 1 plier, the remaining items are 1 plier, 1 nail, 1 bucket. The other player} \\
            }
        }
    }
     \\ \\

    \theutterance \stepcounter{utterance}  
    & & & \multicolumn{2}{p{0.3\linewidth}}{
        \cellcolor[rgb]{0.9,0.9,0.9}{
            \makecell[{{p{\linewidth}}}]{
                \texttt{\tiny{[GM$|$GM]}}
                \texttt{continue} \\
            }
        }
    }
    & & \\ \\

    \theutterance \stepcounter{utterance}  
    & & & \multicolumn{2}{p{0.3\linewidth}}{
        \cellcolor[rgb]{0.9,0.9,0.9}{
            \makecell[{{p{\linewidth}}}]{
                \texttt{\tiny{[GM$|$GM]}}
                \texttt{player was instructed to make a proposal} \\
            }
        }
    }
    & & \\ \\

\end{supertabular}
}

\end{document}
